\subsection*{Axioms}

\begin{remark*}[Source note]
Euclid calls the following principles \textit{common notions}. They function as
general axioms governing equality and comparison in Book I of the
\textit{Elements} \citep[Book I, Common Notions 1--5]{euclid-elements-heath}.
\end{remark*}


\begin{axiombox}{Axiom}
\begin{axiom}[Euclid Common Notion I: Equality is transitive through a common equal]
\label{ax:euclid-common-notion-things-equal-to-same}
Things which are equal to the same thing are also equal to one another.
\end{axiom}
\end{axiombox}
\begin{remark*}[Standard quantified statement]
This is a classical Euclidean source-text clause. In this chapter it is treated
as a named primitive statement rather than expanded into a modern first-order
formalization.
\end{remark*}

\begin{lra-not-visible}
\NoLocalDependencies
\end{lra-not-visible}

\begin{axiombox}{Axiom}
\begin{axiom}[Euclid Common Notion II: Adding equals to equals]
\label{ax:euclid-common-notion-add-equals}
If equals be added to equals, the wholes are equal.
\end{axiom}
\end{axiombox}
\begin{remark*}[Standard quantified statement]
This is a classical Euclidean source-text clause. In this chapter it is treated
as a named primitive statement rather than expanded into a modern first-order
formalization.
\end{remark*}

\begin{lra-not-visible}
\begin{dependencies}
\begin{itemize}
  \item \hyperref[ax:euclid-common-notion-things-equal-to-same]{Euclid Common Notion I}
\end{itemize}
\end{dependencies}
\end{lra-not-visible}

\begin{axiombox}{Axiom}
\begin{axiom}[Euclid Common Notion III: Subtracting equals from equals]
\label{ax:euclid-common-notion-subtract-equals}
If equals be subtracted from equals, the remainders are equal.
\end{axiom}
\end{axiombox}
\begin{remark*}[Standard quantified statement]
This is a classical Euclidean source-text clause. In this chapter it is treated
as a named primitive statement rather than expanded into a modern first-order
formalization.
\end{remark*}

\begin{lra-not-visible}
\begin{dependencies}
\begin{itemize}
  \item \hyperref[ax:euclid-common-notion-things-equal-to-same]{Euclid Common Notion I}
\end{itemize}
\end{dependencies}
\end{lra-not-visible}

\begin{axiombox}{Axiom}
\begin{axiom}[Euclid Common Notion IV: Coincidence implies equality]
\label{ax:euclid-common-notion-coinciding-things-equal}
Things which coincide with one another are equal to one another.
\end{axiom}
\end{axiombox}
\begin{remark*}[Standard quantified statement]
This is a classical Euclidean source-text clause. In this chapter it is treated
as a named primitive statement rather than expanded into a modern first-order
formalization.
\end{remark*}

\begin{lra-not-visible}
\NoLocalDependencies
\end{lra-not-visible}

\begin{axiombox}{Axiom}
\begin{axiom}[Euclid Common Notion V: Whole and part]
\label{ax:euclid-common-notion-whole-greater-than-part}
The whole is greater than the part.
\end{axiom}
\end{axiombox}
\begin{remark*}[Standard quantified statement]
This is a classical Euclidean source-text clause. In this chapter it is treated
as a named primitive statement rather than expanded into a modern first-order
formalization.
\end{remark*}

\begin{lra-not-visible}
\NoLocalDependencies
\end{lra-not-visible}

\begin{remark*}[Exposition]
Common notions are general principles of reasoning rather than specifically
geometric assumptions. They express basic behavior of equality, comparison, and
logical transfer: equal things may be substituted for equal things, adding or
subtracting equals preserves equality, and a whole exceeds its part.

Because these principles are not tied to points, lines, circles, or planes,
they can be used throughout mathematics, not only in geometry. In modern
language, they play the role of background axioms governing how equality and
comparison may be used inside proofs.
\end{remark*}
